% concatenated from S-Noetherian_lattices-_22April_2026.tex
\documentclass[10pt,twoside,a4]{amsart}
\usepackage{amssymb,amsmath,amsfonts,amsthm,graphicx,amscd}
\usepackage{enumerate,mathrsfs,longtable,bm,float,xcolor,tikz,tikz-cd}
\usepackage[dvips]{epsfig}
\usepackage{amsmath,amsfonts,amssymb,amsthm,bm,float,makecell,booktabs,MnSymbol}
%\usepackage{slashbox}

%to check all references are cited or not

%\usepackage{refcheck}
\usepackage{enumerate}
\usepackage{amscd}

\usetikzlibrary{positioning}
\usepackage[all,cmtip,line]{xy}
%\voffset-0.8in \hoffset-.2in
\renewcommand{\baselinestretch}{1.37}
\setlength{\oddsidemargin}{1 cm} \setlength{\evensidemargin}{1cm}
\setlength{\topmargin}{1 cm} \textwidth = 15.8 cm \textheight = 20cm 

\flushbottom

\DeclareMathAlphabet{\mathpzc}{OT1}{pzc}{m}{it}

\numberwithin{equation}{section}

%\hbadness3000 \vbadness30000
%\parindent= 18 pt
\theoremstyle{plain}
\newtheorem{lem}{Lemma}[section]
\newtheorem{Lemma}[lem]{Lemma}
\newtheorem{Theorem}[lem]{Theorem}
\newtheorem{thm}[lem]{Theorem}
\newtheorem{prop}[lem]{Proposition}
\newtheorem{obs}[lem]{Observation}
\newtheorem{cor}[lem]{Corollary}
\theoremstyle{definition}
\newtheorem{exa}[lem]{Example}
\newtheorem{rem}[lem]{Remark}
%\newtheorem{obs}[lem]{Observation}
\newtheorem{defn}[lem]{Definition}
\newtheorem{definition}[lem]{Definition}
\newtheorem{Conj}[lem]{Conjecture}
\newtheorem{que}[lem]{Question}

\newtheorem{prfname}{\it{Proof of Theorem }}
%\renewcommand\labelenumi{\it(\arabic{enumi})}
\newcommand{\Z}{{\mathbb{Z}}}
\newcommand{\zg}{$\Gamma_{I}(P)$ }
\thispagestyle{empty}
\setcounter{page}{1}
\begin{document}
	
	\baselineskip 13truept
	
	\title{On $S$-Noetherian Lattices}
	
	\author{Sachin Sarode*,  Chetan Patil** and Vinayak Joshi***}
	\address{\rm *Department of Mathematics, Shri Muktanand College\\ Gangapur, Dist. Chh. Sambhajinagar - 431 109, India.} \email{sarodemaths@gmail.com}
	
	
	\address{\rm **School of Technology Management and Engineering, SVKM NMIMS Global University,
		Dhule-424 001, India.}
	\email{patilcs19@gmail.com}
	\address{\rm ***Department of Mathematics, Savitribai Phule Pune University, Pune-411 007, India.}
	\email{vvjoshi@unipune.ac.in \\
		vinayakjoshi111@yahoo.com }
	
	
	
	%\thanks{ The second author  gratefully acknowledges the financial support by University Grants Commission, New Delhi, via research project reference number 47-496/12(WRO)}
	\subjclass[2020]{Primary 06F10, 13A15, 13C05,  Secondary 06A11}
	%\today{} 
	
	\maketitle
	

	\begin{abstract} 
 In this paper, we define and study $S$-Noetherian lattices as a natural generalization of Noetherian rings. We prove that a ring $R$ is $S$-Noetherian if and only if its ideal lattice, $Id(R)$, is $S_L$-Noetherian. Furthermore, we establish a Cohen–Kaplansky type theorem for $S$-Noetherian lattices, showing that $L$ is $S$-Noetherian if and only if every $S$-prime element of $L$ is $S$-compact. Finally, we introduce the concept of $S$-primary elements—a generalization of primary elements in multiplicative lattices—and demonstrate the existence and uniqueness of $S$-primary decomposition in $S$-Noetherian lattices.		
	\end{abstract}
	
	\vskip 5truept
	
	\noindent\textbf{Keywords:} $S$-prime element,  $S$-Noetherian lattice, $S$-primary element, multiplicative lattice.
	\vskip 5truept
	
	%\noindent\textbf{Mathematics Subject Classification:}  
	
	
	
	\baselineskip 14truept 
	\section{Introduction}\label{intro}

The study of Noetherian rings remains a cornerstone of commutative algebra, primarily due to  the existence of primary decomposition. In recent years, significant research has focused on generalizing these classical notions to various $S$-variants. A particularly fruitful direction has been the introduction of the $S$-Noetherian property. Originally introduced by Ahmed and Sana \cite{HH} to characterize rings that satisfy a weakened version of the Noetherian condition relative to a multiplicatively closed set $S$, this concept has proven essential for understanding rings that fail the classical $ACC$ but remain manageable under $S$-localized conditions. 

In this paper, we extend this theory to  multiplicative lattices, providing a complete abstraction of $S$-Noetherian rings. Multiplicative lattices serve as an ideal framework for this unification, introduced by Ward and Dilworth \cite{WD}, as they allow for the study of ideal-theoretic properties independent of the underlying element-wise ring structure.  One of our central results, an analogue of Cohen-Kaplansky type theorem,  is a characterization theorem which states that an $r$-lattice $L$ is $S$-Noetherian if and only if every prime element is $S$-compact, which is further shown to be equivalent to $L$ satisfying the $S$-stationary property. In this context, the $S$-stationary property serves as the natural lattice-theoretic $S$-variant of the $ACC$ in rings, while $S$-compact elements is an  abstraction of $S$-finitely generated ideals. 

Massaoud \cite{M} and Visweswaran \cite{V} independently introduced the concept of $S$-primary ideals in commutative rings as a generalization of classical primary ideals. Subsequently, T. Singh, A. Ansari, and S. Kumar \cite{SAK} established the existence of $S$-primary decomposition for $S$-Noetherian rings. 

In the present work, we extend these results to the lattice-theoretic setting. We introduce the concept of $S$-primary elements in multiplicative lattices and establish both the existence and uniqueness of $S$-primary decomposition within an $S$-Noetherian lattice. Finally, we bridge the gap between ring-theoretic results and our lattice-theoretic abstractions. We establish a rigorous correspondence: if $R$ is a commutative ring with identity and $S$ is a multiplicatively closed subset of $R$, then an $S$-primary ideal of $R$ corresponds precisely to an $S_L$-primary element of the ideal lattice $L = Id(R)$, where $S_L$ is the induced multiplicatively closed subset of $L$. This correspondence ensures that our results for $S$-Noetherian lattices directly recover, unify, and extend existing theorems in commutative algebra. 
	
	We begin with the definitions and terminologies.
	
		\begin{defn}
		A \textit{multiplicative lattice} is a complete lattice $L$ together with a binary operation (multiplication) ``$\cdot$" satisfying:
		\begin{enumerate}
			\item $(L,\;\cdot)$ is commutative and associative,
			\item multiplication distributes over arbitrary joins,
			\item $1$ is the multiplicative identity.
		\end{enumerate}
	\end{defn}
	

	

	
	
	


	If $L$ is a multiplicative lattice and $a, b \in L$, then $a\cdot b \leq a \wedge b$. This property follows from the fact that the multiplication distributes over join and 1 is the multiplicative identity. 
	
It is well known that the ideal lattice of a commutative ring with unity is a multiplicative lattice; however, every multiplicative lattice is not an ideal lattice of some commutative ring with unity. For instance, consider the multiplicative lattice $L$ described in Example 2.4, \cite{JS}. Since L is non-modular, it cannot be the ideal lattice of any ring, because the ideal lattice of a ring is necessarily modular.
	

	\begin{defn}[Sarode, Joshi and Patil \cite{SJP}]\label{def:vlattice}
		A complete lattice $L$ is said to be a \textit{$V$-lattice} if there exists a commutative, associative binary operation ``$\cdot$" on $L$ such that:
		\begin{enumerate}
		\item $a \leq b \implies a \cdot c \leq b \cdot c$, for  $a, b, c\in L$.
			\item $a \cdot b \leq a \wedge b$, for  $a, b\in L$.
			\item $a \cdot 1 = a$ for every $a\in L$.
		\end{enumerate}
	\end{defn}
	
	\begin{rem} \label{1.7.}
		Every multiplicative lattice is a $V$-lattice, but not conversely.
	\end{rem}


		
		\begin{defn}
			A nonempty subset $S$ of $L_{*}$ in a $V$-lattice $L$ is \textit{multiplicatively closed} if $s_{1} \cdot s_{2} \in S$ whenever $s_{1}, s_{2} \in S$. Here    $L_{*}$ denotes the set of all compact elements of a complete lattice $ L $. By a \textit{compact} element $c$ of a complete lattice $ L $, we mean 
			if $ c \leq \bigvee_{\alpha} a_{\alpha} $, $\alpha \in \Lambda$
			($\Lambda$ is an indexed set) implies $ c \leq \bigvee _{i=1}^{n}
			a_{\alpha_{i}} $, where $ n \in \mathbb{Z}^{+} $. A lattice $L$ is \textit{compactly generated} if every
			element of $L$ is a join of compact elements of $L$.  
				\end{defn}
			\begin{defn} [ \cite{AAJ, A, SJ, WD}] A multiplicative lattice $L$ is \textit{$1$-compact} if $1$ is a compact element of $L$. A multiplicative lattice $L$ is a {\it $c$-lattice} if $L$ is 1-compact, compactly generated multiplicative lattice in which the product of two compact elements is compact. \end{defn}	
	
		

		 \textbf{Throughout this paper, $S$ is a multiplicatively closed subset of a $V$-lattice $L$ such that $1 \in S$ and $0 \notin S$.}
		
		
		\begin{defn}[Sarode, Joshi and Patil \cite{SJP}]
		Let $S$ be a multiplicatively closed subset of a $V$-lattice $L$. A proper element $p$ of $L$ such that $t \not\leqq p$ for all $t \in S$. Then $p$ is an { \it $S$-prime} element of $L$ if there exists an element $s \in S$ such that for all $a, b  \in L$ with $a \cdot b \leq p$, implies $s \cdot a \leq p$ or $s \cdot b \leq p$. 
			
		If we take $S=\{1\}$, then we get the classical definition of prime elements. That is, 	an element $p \neq 1$ of a $V$-lattice $L$  is a \textit{prime} element, if $a\cdot b \leq p$ implies $a \leq p$ or $b \leq p$ for all $a, ~b \in L$.  
		
		 Equivalently, in a $c$-lattice $L$, an element $p \neq 1$  is a \textit{prime} element, if $a\cdot b \leq p$ implies $a \leq p$ or $b \leq p$ for  $a, ~b \in L_{\star}$, the set of all compact elements.

		The set of all $S$-prime elements of $L$ is denoted by $Spec_s(L)$. If $S = \{1\}$, then $Spec_s(L)$ is denoted by $Spec(L)$, the set of all prime elements of $L$.
	\end{defn}
		
	%\begin{exa}\label{34.6.} 	Consider the lattice $N_5$, whose Hasse diagram is shown in Figure \ref{fig:n5}.   Note that $N_5$ is a $V$-lattice with the multiplication as meet. If we take $S = \{1 \}$, then $b$ and $c$ are $S$-prime elements of $N_5$, however, $a$ and $0$ are not $S$-prime elements of $N_5$. If  $S = \{1, c \}$, then $0, ~ a$ and $b$ are $S$-prime elements of $N_5$, however, $c$ is not an $S$-prime element of $N_5$, as $c\leq c\in S$.  Further  if $S = \{1, ~a \}$, then $0$ and $c$ are $S$-prime elements of $N_5$ but $a$ and $b$ are not  $S$-prime elements of $N_5$. \end{exa}
		
		
	%\begin{exa}\label{44.6.}	Consider the lattice $L=Id(\mathbb{Z}_{12})$, the ideal lattice of the ring $\mathbb{Z}_{12} $.  Observe that  $L$ is a multiplicative lattice with the multiplication as the multiplication  of ideals. If we take $S = \{(1), ~(4) \}$, then $(0), ~ (6)$, and $ (3)$ are $S$-prime elements of $L$, however, $(2)$ and $(4)$ are not $S$-prime elements of $L$. If $S = \{(1), ~(3) \}$, then $(2)$ and $ (6)$ are $S$-prime but $(0), ~(3)$, and $(4)$ are not $S$-prime.    \end{exa}			
			
		

	%	{\color{red} \begin{rem} 			$S$-prime elements defined in this paper is different than $s$-prime elements defined in Jayram and Johnson \cite{JJ}. 		\end{rem}}
		
		
	Following result is due to Hamed and Malek \cite{HM}.
		
	\begin{thm}[Hamed and Malek \cite{HM}] \label{1.20.}
	Let $R$ be a commutative ring with identity, $S $ be a multiplicative subset of $R$ and $P$ be an ideal of $R$ disjoint with $S$. Then $P$ is $S$-prime if and only if there exists $s \in S$, such that for all $I, J$ two ideals of $R$, if $IJ \subseteq P$, then $sI \subseteq P$ or $sJ \subseteq P$.
	\end{thm}
		
		
	\begin{thm}[Sarode, Joshi and Patil \cite{SJP}]
	Let $R$ be a commutative ring with identity and $S$ a multiplicative subset of $R$. Let $L=Id(R)$ denote the multiplicative lattice of all ideals of $R$, and define 
	$S_L=\{(a)\in L \mid a\in S\}.$ Then an ideal $P$ of $R$ is an $S$-prime ideal of $R$ if and only if $P$ (viewed as an element of $L$) is an $S_L$-prime element of $L$.
	\end{thm}
		
	
	Thus the lattice-theoretic notion of $S$-prime elements in $L$  extends the ring-theoretic notion of $S$-prime ideals in $R$.
		
		
		\begin{defn} [ \cite{D, SJ, WD}]
			Let $a,b $ be any elements of    a multiplicative lattice $L$. Then $(a:b) =
			\bigvee \{x \;|\;  x\cdot b \leq a\}$. Note that $x \cdot b \leq a
			\Leftrightarrow x \leq (a:b)$.  Let $L$ be a $c$-lattice and $a\in L$. The \textit{radical} of $a$ is denoted by $\sqrt a$ is defined as $\sqrt a =
			\bigvee \{x\in L_* \;|\;  x^n \leq a ~ \textrm{for some }~ n\in \mathbb{Z}^+\}$. \end{defn}
		Using the fact that the compact elements are precisely the finitely generated ideals.  We have:
		
		\begin{lem}\label{lem:radical-equals-ring}
			Let $R$ be a commutative ring with $1$, and let $L=Id(R)$ be the multiplicative lattice of ideals of $R$. For $A\in L$,  the lattice-theoretic radical
			$	\sqrt{A}	=\bigvee\{\,X\in L_* \mid X^n \leq A \text{ for some } n\in\mathbb{Z}^+ \}$ coincides with $\sqrt{A} = \{\, r\in R \mid r^n \in A \text{ for some } n\in\mathbb{Z}^+ \}$,  i.e., the lattice-theoretic radical coincides with the usual radical ideal of $A$ in $R$.
		\end{lem}
		
		
%\begin{defn} [ \cite{AAJ, A, SJ, WD}] A multiplicative lattice $L$ is \textit{$1$-compact} if $1$ is a compact element of $L$. A multiplicative lattice $L$ is a {\it $c$-lattice} if $L$ is 1-compact, compactly generated multiplicative lattice in which the product of two compact elements is compact. \end{defn}
		



\section{S-Noetherian Lattices}


Several authors, including Ward and Dilworth \cite{WD}, have studied \textit{multiplicative lattices} in connection with the ideal theory of rings. To generalize the study  of Noetherian rings to a broader order-theoretic framework, Jayaram\cite{JG} studied the concept of a \textit{Noetherian  lattice}. It is well known that the ideal lattice $Id(R)$ of a commutative ring $R$ forms a compactly generated, modular lattice in which the greatest element $1$ is compact. In fact, the notion of a Noetherian lattice is a true generalization of that of a Noetherian ring, as evidenced by the equivalence
$Id(R) $ is a Noetherian lattice  if and only if $R$  is a Noetherian ring.
This correspondence establishes a strong structural bridge between ring-theoretic and lattice-theoretic finiteness conditions.

In this section, we introduce an \textit{$S$-version of Noetherian lattices}, extending the classical framework by incorporating a multiplicatively closed subset $S$ of a $V$-lattice $L$. Motivated by the notion of  \linebreak $S$-Noetherian rings in commutative algebra, we define an \textit{$S$-Noetherian multiplicative lattice} and an analogue of the ascending chain condition (ACC) in terms of the \textit{$S$-stationary property}. Thereby providing a flexible generalization  within the lattice context.

The study of $S$-Noetherian  lattices therefore lies at the intersection of \textit{order theory} and the \textit{ideal theory of commutative rings}, serving as a unifying framework to generalize classical algebraic concepts. This perspective offers an elegant order-theoretic setting to investigate extensions of Noetherian properties, finiteness conditions, and chain conditions beyond the algebraic confines of ring theory.



\begin{defn}[D. D. Anderson \cite{A}]
	Let $a, b, m$ be elements of a multiplicative lattice $L$. 
	\begin{enumerate}
		\item An element $m$ is said to be\textit{ meet principal} if $a \wedge mb = m((a:m) \wedge b)$ for all $a,b \in L$.
		\item An element $m$ is said to be \textit{join principal} if $a \vee (b :m) = (am \vee b): m$ for all $a,b \in L$.
		\item An element $m$ is said to be  \textit{principal} if $m$ is both meet principal and join principal.
	\end{enumerate}
	
\end{defn}

\begin{defn} [D. D. Anderson \cite{A}]
	A multiplicaive lattice $L$ is said to be an \textit{$r$-lattice}, if  $L$ is a modular, principally generated, and compactly generated lattice with $1$ as a compact element of $L$. 
\end{defn}

The following definition of $S$-Noetherian ring is due to Anderson and Dumitrescu \cite{AT}. 

\begin{defn} [D. D. Aderson and Tiberiu Dumitrescu \cite{AT}]
	Let $S$ be a multiplicative subset of a commutative ring $R$ with 1. An ideal $I$ of $R$ is said to be \textit{$S$-finite}, if there exists some finitely generated ideal $J$ of $R$ and some $s \in S$ such that $sI \subseteq J \subseteq I$. A commutative ring $R$ with 1 is said to be \textit{$S$-Noetherian} if each ideal of $R$ is $S$-finite.  
\end{defn}

 
Now, analogously we define  $S$-Noetherian lattice as follows.

\begin{defn}
	Let $S$ be a multiplicative subset of a multiplicative lattice $L$. An element $a$ of $L$ is said to be \textit{$S$-compact}, if there exists some compact element $b$ in $L$ and some $s \in S$ such that $s\cdot a \leq b \leq a$. A $r$-lattice $L$ is said to be an \textit{$S$-Noetherian lattice} if each element of $L$ is $S$-compact.  
\end{defn}



In the following result, we connect the $S$-Noetherian rings with $S_L$-Noetherian lattices.

\begin{thm}\label{thm:SL-Noeth}
	Let $R$ be a commutative ring with $1$ and let $S\subseteq R$ be a multiplicatively closed set.
	Set $L:=Id(R)$ and $S_L:=\{\, (s)\in L \mid s\in S \,\}$.
	Then the following are equivalent:
	\begin{enumerate}
		\item $R$ is $S$-Noetherian (i.e., every ideal $I$ is $S$-finite).
		\item $L=Id(R)$ is $S_L$-Noetherian as a multiplicative lattice (i.e., every element $I\in L$ is $S_L$-compact).
	\end{enumerate}
\end{thm}

\begin{proof}
	Note that $Id(R)$ is an $r$-lattice.
	
	$(1)\Rightarrow(2)$:
	Let $I\in L$. As $R$ is $S$-Noetherian, there exist $s\in S$ and a finitely generated ideal $J\subseteq I$ with $sI\subseteq J$.
	In $L$, $(s)\cdot I  \leq J \leq I$, and $J$ is compact. Hence $I$ is $S_L$-compact. As $I$ was arbitrary, $L$ is $S_L$-Noetherian.
	
	$(2)\Rightarrow(1)$:
	Let $I$ be an ideal of $R$. Since $L$ is  $S_L$-Noetherian, there exists $(s)\in S_L$ and a compact $b\in L$ with $(s)\cdot I\leq b\leq I$.
	It is well known that the compact elements in $Id(R)$ are precisely the  finitely generated ideals of $R$, we get $sI\subseteq b\subseteq I$ with $b$ as a finitely generated ideal. Thus $I$ is $S$-finite. Since $I$ was arbitrary, $R$ is $S$-Noetherian.
\end{proof}
\begin{rem}\label{remark}
	Let $L$ be $r$-lattice. Then every principal element of $L$ is compact, and moreover the product of two compact elements is compact in $L$. For more details; see \cite{A}. 
\end{rem}

\begin{lem} \label{5.51.}
	Let $S$ be a multiplicatively closed subset of a $r$-lattice $L$. Then the following are true.
	\begin{enumerate}
		\item If $j$ is a principal element and $i$ is an $S$-compact element of $L$, then $i\cdot j$ is $S$-compact.
		\item If $k$ is a compact element and $i$ is an $S$-compact element of $L$, then $k \vee i$ is $S$-compact.
	\end{enumerate}
\end{lem}

\begin{proof}
	\begin{enumerate}
		\item Since $i$ is an $S$-compact element of $L$, there exist a compact element $b$ and some $s \in S$ such that $s \cdot i \leq b \leq i$. Therefore we have, $s\cdot i \cdot j \leq b\cdot j \leq i\cdot j$. As $j$ is a principal element and $b$ is a compact element of $L$, then by Remark \ref{remark}, $j$ is a compact element of $L$. Moreover $b\cdot j$ is a compact element of $L$.  Therefore 	$i\cdot j$ is an $S$-compact element of $L$.	
		
		\item Since $i$ is an $S$-compact element of $L$, there exist a compact element $b$ and some $s \in S$ such that $s \cdot i \leq b \leq i$. Therefore we have, $(k \vee s i) \leq ( k \vee b) \leq (k \vee i)$. This gives  $s\cdot (k \vee i) \leq  (k \vee s \cdot i) \leq (k \vee b) \leq (k \vee i)$.  Thus $k \vee i$ is an $S$-compact element of $L$.\end{enumerate}\end{proof}




\begin{definition}[Jayaram\cite{JG}]
	A multiplicative lattice $L$ is a Noetherian  lattice, if $L$ is a	$r$-lattice with ACC. Equivalently, $L$ is Noetherian if and only if every element is compact.
\end{definition}

As established earlier, the ideal lattice $Id(R)$ of a Noetherian (or $S$-Noetherian) ring $R$ posses Noetherian (or $S$-Noetherian) property. This correlation naturally invites an investigation into the specific distinction between $S$-Noetherian and classical Noetherian lattices. In this context, a foundational example provided in \cite[Example 2.5]{SAK} serves the purpose. Specifically, consider the following example.

\begin{exa}[\cite{SAK}]Let $L = \mathrm{Id}(R)$ denote the ideal lattice of the polynomial ring $R = F[x_1, x_2, \dots, x_n, \dots]$ over a field $F$. It is well-known that the lattice $L$ is not Noetherian, as it admits the strictly ascending infinite chain of ideals $(x_1) \subsetneqq (x_1, x_2) \subsetneqq (x_1, x_2, x_3) \subsetneqq \dots $ However, $R$ is an $S$-Noetherian ring for the multiplicative set $S = R \setminus \{0\}$. Consequently, by applying Theorem \ref{thm:SL-Noeth}, it follows that the lattice $L$ is $S_L$-Noetherian, where $S_L = \{ (s) \mid s \in S \}$.\end{exa}



Ahmed and Sana \cite{HH} has defined $S$-stationary for sequence of submodules of a module. Analogously, we define $S$-stationary for sequence of elements of a $V$-lattice $L$.

\begin{defn}
	Let $S$ be a multiplicatively closed subset of a $V$-lattice $L$. An ascending chain of elements $i_1 \leq i_2 \leq i_3 \leq \cdots \leq i_m  \leq \cdots $ is called \textit{$S$-stationary} if there exist $s \in S$ and $ n \in \mathbb{Z}^{+}$ such that $s  i_m \leq i_n$ for all $m \geq n$. 


Let $S$ be a multiplicatively closed subset of a $V$-lattice $L$, then  $L$ is said to be satisfies \textit{$S$-stationary property}, if every ascending chain of elements is $S$-stationary.

	Let $S$ be a multiplicatively closed subset of a $c$-lattice $L$ and let $T\subseteq L$. An element $m$ of $T$ is called \textit{$S$-maximal element of $T$} if there exists an element $s\in S$ such that for each $a\in T$, if $m\leq a$, then $sa\leq m$.	
\end{defn}

\begin{thm} \label{5.11a.}
	Let $L$ be a $r$-lattice. Then the following statements are equivalent:
	
	\begin{enumerate}
		\item $L$ is $S$-Noetherian lattice.
		
		\item Every element of $L$ is $S$-compact. 
		
		\item $L$  satisfies $S$-stationary property. 
		\item\label{S-MC} Every non-empty subset of $L$ has $S$-maximal element.
	\end{enumerate}
\end{thm}

\begin{proof}
	$(1) \Leftrightarrow  (2)$ follows from the definition of an $S$-Noetherian lattice.
	
	
	%$(3) \Rightarrow (2)$ Suppose that $L$  satisfies $S$-stationary property. Let $i$ be any element of $L$. Since $L$ is compactly generated  $i = \bigvee_{\alpha} a_{\alpha}$, $\alpha \in \Lambda$ ($\Lambda$ is an index set) and $a_\alpha \in L_*$. We have ascending sequence of elements $a_1 \leq a_1 \vee a_2 \leq a_1 \vee a_2 \vee a_3 \leq \cdots$ in $L$. Since $L$ satisfies $S$-stationary property. We have $s \cdot (a_1 \vee a_2 \vee \cdots \vee a_n \vee a_{n+1} \vee \cdots \vee a_m \vee \cdots) \leq (a_1 \vee a_2 \vee \cdots \vee a_n)$ for all $m \geq n$. Since finite join of compact element is compact, we have $a_1 \vee a_2 \vee \cdots \vee a_n$ is compact element.  Therefore $s \cdot i \leq a_1 \vee a_2 \vee \cdots \vee a_n \leq i$. This proves  that $i$ is $S$-compact.
	
	
	$(2) \Rightarrow (3)$ Suppose that every element of $L$ is $S$-compact. We will show that $L$ satisfies $S$-stationary property. Let $b_1 \leq b_2 \leq b_3 \leq \cdots$ be any ascending chain of elements in $L$. Let $j = \bigvee_{\alpha} b_{\alpha}$, $\alpha \in \Lambda$ ($\Lambda$ is an index set). Since $j$ is $S$-compact, we have $s \cdot j \leq k \leq j$ for some compact element $k$ and some $s \in S$. Therefore $s \cdot (b_1 \vee b_2 \vee b_3 \vee \cdots) \leq k \leq (b_1 \vee b_2 \vee b_3 \vee \cdots)$. Since $k$ is compact and $k \leq (b_1 \vee b_2 \vee b_3 \vee \cdots)$, we have $k \leq b_1 \vee b_2 \vee b_3 \vee \cdots \vee b_n =b_n$, as $b_1 \leq b_2 \leq b_3 \leq \cdots \leq b_n$. So $s \cdot (b_1 \vee b_2 \vee b_3 \vee \cdots) \leq b_n \leq (b_1 \vee b_2 \vee b_3 \vee \cdots)$. This gives $s \cdot b_1 \vee s \cdot b_2 \vee s \cdot b_3 \vee \cdots \leq  b_n$. Thus,  $s b_m \leq b_n$ for all $m \geq n$. Hence $L$  satisfies $S$-stationary property. 
	
	
$(3) \Rightarrow (4)$: Let $T$ be a non-empty subset of $L$. Let $x_1\in T$. If $x_1$ is not $S$-maximal in $T$, then there exist an element $x_2\in T$ such that $x_1\leq x_2$ and $sx_2\nleq x_1$ for all $s\in S$. Again, if $x_2$ is not $S$-maximal in $T$, then there exist an element $x_3\in T$ such that $x_2\leq x_3$ and $sx_3\nleq x_2$ for all $s\in S$. Continuing in this way, we get an ascending chain $x_1\leq x_2\leq x_3\leq \dots$ of elements in $T$ such that $sx_m\nleq x_n$ for all $m\geq n$, $m,n\in \mathbb{Z_+}$ and for all $s\in S$, a contradiction to $(3)$. Therefore, $T$ has an $S$-maximal element. Thus, every non-empty subset of $L$ has an $S$-maximal element.
	
	$(4) \Rightarrow  (2)$: Assume that every non-empty subset of $L$ has an $S$-maximal element. Let $x\in L$.
	
	Since $L$ is an $r$-lattice, it is compactly generated and hence there exists a family $\{a_\alpha\}_{\alpha\in\Lambda}\subseteq L_*$
	of compact elements such that
	$x=\bigvee_{\alpha\in\Lambda} a_\alpha$. Let
	$T=\left\{\bigvee_{i=1}^n a_{\alpha_i}\mid
	n\in\mathbb{N},\ \alpha_1,\dots,\alpha_n\in\Lambda\right\}$.
	Then $T$ is a non-empty subset of $L$, and every element of $T$ is compact,
	since finite join of compact elements is compact.
	Moreover, $\bigvee T=x$.
	
	By the assumption, $T$ has an $S$-maximal element; say, $k\in T$ be such an element.
	Then $k$ is compact and $k\leq x$. Since $k$ is $S$-maximal in $T$ and $k\leq x$, there exists $s\in S$ such that $sx\leq k$.
	 Therefore, $sx\leq k\leq x$, with $k$ compact.
	Hence $x$ is an $S$-compact element of $L$. Since $x\in L$ was arbitrary, every element of $L$ is $S$-compact.
	\end{proof}


In sequel, we obtain the S-version of Cohen's result. For details of Cohen's result; see Kaplansky \cite[Theorem 8]{kap}.  

\begin{rem}\label{princ-prime}
	\begin{enumerate}
	\item Observe that if $L$ is an $r$-lattice, then to establish that an element 
	$p \in L$ is prime, it is enough to verify that
	$x \cdot y \leq p \;\implies\; x \leq p \ \text{or}\ y \leq p
	$
	holds whenever $x, y$ are principal elements of $L$.
	
	\item Note that if $a$ is not an $S$-compact element of $L$, then $t\nleq a$ for all $t\in S$. For, if $t_1\leq a$ for some $t_1\in S$, then $t_1a\leq t_1\leq a$, where $t_1\in S\subseteq L_{\star}$, the set of compact elements. Therefore, $a$ is  $S$-compact, a contradiction. 
	\end{enumerate}
\end{rem}




\begin{thm}[Cohen--Kaplansky $S$-version]
	Let $S$ be a multiplicatively closed subset of an $r$-lattice $L$. Then $L$ is an $S$-Noetherian lattice if and only if  every $S$-prime element of $L$  is  $S$-compact. 
\end{thm} 


\begin{proof}If $L$ is an $S$-Noetherian lattice, then every element is $S$-compact and hence every prime element is $S$-compact. 
	
	
	Conversely, assume that every $S$-prime element of $L$  is $S$-compact. Let $\beta = \{a' \in L\; |\; a'$ is not  $S$-compact$\}$. If $\beta = \emptyset$, then all elements of $L$ are $S$-compact and by Theorem \ref{5.11a.}, $L$ is an $S$-Noetherian lattice.
	
	Now, assume that $\beta \neq \emptyset$. Let $a_{\alpha} \in \beta$, $\alpha \in \Lambda$ ($\Lambda$ is an index set). Let $ a_1 \leq a_2 \leq a_3 \leq \cdots$ be increasing sequence of elements from $\beta$. Let $a = \bigvee_{\alpha} a_{\alpha}$, $\alpha \in \Lambda$ ($\Lambda$ is an index set). We show that $a$ is upper bound of $a_1 \leq a_2 \leq a_3 \leq \cdots$ and $a \in \beta$.
	
	First, we claim that $a$ is not an $S$-compact element of $L$. Suppose on the contrary that $a$ is an $S$-compact element of $L$. Therefore there exists $s \in S$ and a compact element $b \in L$ such that $s \cdot  a \leq b \leq a$, that is $s \cdot (\bigvee_{\alpha} a_{\alpha}) \leq b \leq (\bigvee_{\alpha} a_{\alpha})$. Since $b$ is compact and  $b \leq (\bigvee_{\alpha} a_{\alpha})$, we have $b \leq a_1 \vee a_2 \vee \cdots \vee a_n$ for some $ n \in \mathbb{Z}^{+}$. As  $a_1 \leq a_2 \leq a_3 \leq \cdots$ is ascending chain of elements and $b \leq a_1 \vee a_2 \vee \cdots \vee a_n=a_n$ for some $n \in \mathbb{N}$,  we have $s \cdot (\bigvee_{\alpha} a_{\alpha})  \leq b \leq a_n$. Since $s \cdot a_n \leq s \cdot (\bigvee_{\alpha} a_{\alpha})$, we have $s \cdot  a_n \leq b\leq a_n$, a contradiction to $a_n$ is not an $S$-compact element. Therefore  $a$ is not an $S$-compact element of $L$.  Hence,  $a \in \beta$. Moreover, $a$ is an upper bound of  $a_1 \leq a_2 \leq a_3 \leq \cdots$. By Zorn's Lemma, $\beta$ contains a maximal element, say $p$. Therefore, $p$ is not an $S$-compact element. This gives $t\nleq p$ for all $t\in S$.
	If $p$ is prime, then $p$ is also $S$-prime for the multiplicatively closed subset $S$ of $L$. Hence, by hypothesis, $p$ is an $S$-compact element, contradicting $p\in\beta$ and in this case we are through. 
	
	If $p$ is not a prime element, then, by Remark \ref{princ-prime}, there exist principal elements $x, y$ in $L$ such that $xy \leq p$ with $x \not\leqq p$ and $y \not\leqq p$. Let $A = x \vee p$ and $B = (p : x)$. Clearly, $p<A,B$, as $x\leq A$, but $x\nleq p$ and $y\leq B$, but $y\nleq p$. Therefore, $A,B\notin \beta$. Hence, $A$ and $B$ are $S$-compact elements of $L$. Since $x$ is principal, $p \wedge x = x \cdot (p: x) = x \cdot  B$. By Lemma \ref{5.51.}, $x \cdot B$ is $S$-compact and hence $p \wedge x$ is $S$-compact. Since $L$ is modular, by  isomorphism theorem for modular lattices, we have 
	$[p\wedge x, p] \cong [x, p\vee x]$ and hence $p = y \vee (p \wedge x)$, where $y$ is a compact element. By Lemma \ref{5.51.}, $p$ is $S$-compact, contradicting $p\in\beta$.
	Thus, $\beta=\varnothing$, i.e., every element of $L$ is $S$-compact, and by Theorem~\ref{5.11a.}, $L$ is $S$-Noetherian.
\end{proof}



\section{$S$-primary decomposition in $S$-Noetherian lattice}	
In this section, we have given the existence and uniqueness of  $S$-primary decomposition in an $S$-Noetherian lattice.  This extends the results of $S$-Noetherian rings to $S$-Noetherian lattices.

\begin{defn}[\cite{M,V}]
	Let $R$ be a commutative ring with identity and $S$ a multiplicatively closed subset of $R$.
	A proper ideal $Q$ of $R$ is called an \emph{$S$-primary ideal} of $R$ if
	$S\cap Q=\varnothing$ and there exists $s\in S$ such that for all $a,b\in R$,
	$ab\in Q$ implies that $sa\in Q$ or $sb\in \sqrt{Q}$.
\end{defn}

Analogously, we define an $S$-primary element in a multiplicative lattices as follows.

\begin{definition}  Let $S$ be a multiplicatively closed subset of a $c$-lattice $L$. A proper element $q$ of $L$ such that $t\nleq q$ for all $t\in S$ is called an \textit{$S$-primary } if there exists an element $s\in S$ such that for all $c,d \in L$ with $cd \leq q$  implies $sc \leq q$ or $sd \leq \sqrt q$. 
	
	
	If $q$ is an $S$-primary element of a $c$-lattice $L$ and  $\sqrt{q}=p$ is $S$-prime element, then  $q$ is called $S$-$p$-primary element of $L$. 
\end{definition}



\begin{rem}
	\begin{enumerate}
		
		\item Let $L$ be a $c$-lattice and $S=\{1\}$, then the primary elements of $L$ and the $S$-primary elements of $L$ coincide.
		\item  If $q$ is a primary element of a $c$-lattice $L$ and $S$ is a multiplicatively closed subset of $L$ such that $t\nleq q$ for all $t\in S$, then $q$ is a $S$-primary element of $L$.
		
	\end{enumerate}
\end{rem}

\begin{exa}	
	\begin{enumerate}
		\item Consider a lattice $L=Id(\mathbb{Z}_{12})$, the ideal lattice of the ring $\mathbb{Z}_{12} $.  Observe that  $L$ is a c-lattice with the multiplication as the multiplication  of ideals. Let $S = \{(1), ~(3) \}$. Then
		\begin{enumerate}
			\item  $(4)$ is an $S$-primary element of $L$ but not an $S$-prime element of $L$.
			
			\item $(6)$ is an $S$-primary element of $L$ but not a primary element of $L$.
		\end{enumerate}
		
		\item Let $L=Id(\mathbb{Z}_{12})$. Let $S = \{(1), ~(4) \}$. Then $(4)$ is a primary element of $L$ but not $S$-primary.
		
		%\item Consider $L=Id(\mathbb{Z}_{12})$. Let $S = \{(1)\}$. Then the primary elements of $L$ and the $S$-primary elements of $L$ coincide.
	\end{enumerate}	
\end{exa}

The following result gives the characterization of $S$-primary ideals.
\begin{thm}[\cite{M}, Theorem 2.11]\label{2.11.}
	Let $R$ be a commutative ring with identity, $S$ a multiplicative subset of $R$ and $P$ an ideal of $R$ disjoint from $S$. Then $P$ is $S$-primary if and only if there exists $s \in S$, such that for all $I, J$ two ideals of $R$, if $IJ \subseteq P$, then $sI \subseteq P$ or $sJ \subseteq \sqrt{P}$.
\end{thm}



\begin{thm}\label{thm:S-primary-lattice-converse}
	Let $R$ be a commutative ring with identity and $S$ a multiplicatively closed subset of $R$. 
	Let $L=Id(R)$ denote the multiplicative lattice of all ideals of $R$, and define
	$S_L=\{(a)\in L \mid a\in S\}$.
	Then an ideal $Q$ of $R$ is an $S$-primary ideal of $R$ if and only if $Q$ (viewed as an element of $L$)
	is an $S_L$-primary element of $L$.
\end{thm}

\begin{proof}
	Suppose that $Q$ is an $S$-primary ideal of $R$.  
	Then $S\cap Q=\varnothing$. This shows that  $t\not\leqq Q$ for all $t\in S_L$.
	Further, for all $I,J\in L$, $IJ\leq Q$ implies that there exists $s\in S$ such that
	$(s)I\leq Q$ or $(s)J\leq \sqrt{Q}$.
	Hence $Q$ is an $S_L$-primary element of $L$.
	
	Conversely, assume that $Q$ is an $S_L$-primary element of $L$.  
	By the definition, $t\not\leqq Q$ for all $t\in S_L$.  
	If there exists $a\in S\cap Q$, then $(a)\in S_L$ and $(a)\leq Q$, a contradiction.
	Thus $S\cap Q=\varnothing$.
	
	Now let $I,J$ be ideals of $R$ such that $IJ\subseteq Q$.
	Equivalently, $IJ\leq Q$ in $L$.
	Since $Q$ is $S_L$-primary, there exists $(s)\in S_L$, for some $s\in S$, such that
	$(s)I\leq Q$ or $(s)J\leq \sqrt{Q}$.
	That is, $sI\subseteq Q$ or $sJ\subseteq \sqrt{Q}$, by Lemma \ref{lem:radical-equals-ring}.
	Hence $Q$ is an $S$-primary ideal of $R$.
\end{proof}

\begin{definition}
	Let $L$ be a multiplicative lattice and $S \subseteq L$ be a multiplicatively closed subset. Let $a \in L$ be an element such that its $S$-saturation $a_S < 1$. We say that $a$ admits an \textit{$S$-primary decomposition } if it can be represented as a finite meet of $S$-primary elements
	$a = \bigwedge_{i=1}^{n} q_i$ where each $q_i$ is an $S$-$p_i$-primary element of $L$. Such a decomposition is said to be \textit{minimal }  if the following two conditions are satisfied:
	\begin{itemize}
		\item The $S$-saturations of the radicals are distinct for all $i, j \in \{1, 2, \dots, n\}$:
		$(p_i)_S \neq (p_j)_S \;\; \text{for } i \neq j$.
		\item   For each $i \in \{1, 2, \dots, n\}$:
		$\left( \bigwedge_{j \neq i} q_j \right)_S \not\leq (q_i)_S$ (Equivalently, for each $i$, there exists an element $x_i \in \bigwedge_{j \neq i} q_j$ such that $s x_i \not\leq q_i$ for all $s \in S$.)
	\end{itemize} 
\end{definition}



\begin{prop}\label{prop:finite-intersection-S-primary}
	Finite meet of $S$-$P$-primary elements is an $S$-$P$-primary element in a $c$-lattice $L$.
\end{prop}

\begin{proof}
	Let $q_1,q_2,\dots,q_n$ be $S$-$P$-primary elements of $L$.
	Then $t\nleq q_i$ for all $t\in S$ and for each $i=1,2,\dots,n$.
	Hence $t\nleq \bigwedge_{i=1}^{n} q_i$ for all $t\in S$. Since $L$ is a $c$-lattice, we have
	$\sqrt{\bigwedge_{i=1}^{n} q_i}
	=\bigwedge_{i=1}^{n}\sqrt{q_i}
	=p$.
	
	Now let $x,y\in L$ such that
	$xy\leq \bigwedge_{i=1}^{n} q_i$
	and suppose that $sy\nleq \bigwedge_{i=1}^{n} q_i$ for all $s\in S$.
	If there exists $s\in S$ such that $sy\leq q_i$ for all $i$,
	then $sy\leq \bigwedge_{i=1}^{n} q_i$, a contradiction.
	Therefore, for each $s\in S$, there exists some index $k$ such that
	$sy\nleq q_k$.
	
	Since $xy\leq \bigwedge_{i=1}^{n} q_i$, we have $xy\leq q_k$.
	As $q_k$ is an $S$-$P$-primary element, there exists $s'\in S$ such that
	$s'x\leq \sqrt{q_k}=p=\sqrt{\bigwedge_{i=1}^{n} q_i}$.
	Thus $\bigwedge_{i=1}^{n} q_i$ is an $S$-$P$-primary element of $L$.
\end{proof}


%\begin{definition}	Let $S$ be a multiplicatively closed subset of a $c$-lattice $L$. An element $m$ of $L$ such that $t\nleq m$ for all $t\in S$ is called \textit{S-maximal element of $L$} if there exists an element $s\in S$ such that whenever $m\leq a$ for some $a\in L$, then either $sa\leq m$ or $s\leq a$.	\end{definition}		

\begin{prop}\label{} 
	Let $S$ be a multiplicatively closed subset of a $c$-lattice $L$. Let $q\in L$ such that $t\nleq q$ for all $t\in S$. If $\sqrt{q}=m$ is $S$-maximal element of $L$, then $q$ is $S$-primary element of $L$.
\end{prop}

\begin{proof} 
	Let $x, y \in L$ such that  $xy\leq q$ with $sy\nleq \sqrt{q}=m$ for all $s\in S$. Let $m^{\prime}\in L$ be a maximal element such that $m\leq m^{\prime}$  and $t\nleq m^{\prime}$ for all $t\in S$. Since $m$ is $S$-maximal element of $L$, there exists $s^{\prime}\in S$ such that $s^{\prime}m^{\prime}\leq m$. Therefore $sy\nleq s^{\prime}m^{\prime}$ for all $s\in S$.  In particular,  $s^{\prime}y\nleq s^{\prime}m^{\prime}$. Also, if $y\leq m^{\prime}$, then $s^{\prime}y\leq s^{\prime}m^{\prime}$, a contradiction. Therefore $y\nleq m^{\prime}$. Hence $m^{\prime} \vee y=1$. Then $s^{\prime}=s^{\prime}(1)=s^{\prime}(m^{\prime} \vee y)=s^{\prime}m^{\prime}\vee s^{\prime}y$. Since $s^{\prime}m^{\prime}\leq m=\sqrt{q}$, there exists $k\in \mathbb{N}$ such that $(s^{\prime})^k(m^{\prime})^k\leq q$. Now, $(s^{\prime})^kx=(s^{\prime}m^{\prime}\vee s^{\prime}y)^k x=(s^{\prime})^k (m^{{\prime}k}\vee m^{{\prime}{k-1}}y\vee \dots \vee m^{\prime}y^{k-1}\vee y^k)x \leq (s^{\prime})^k(m^{{\prime}k} \vee y)x= (s^{\prime})^k(m^{\prime})^k x\vee (s^{\prime})^kyx\leq q\vee q=q $. Hence $ (s^{\prime})^kx\leq q$.\end{proof}	

Now, we relate the known  concept of $S$-irreducible ideals in a commutative ring with the $S$-irreducible elements in a multiplicative lattices.

\begin{definition} (Definition 2.2 \cite{SAK})
	Let $R$ be a commutative ring with identity and let $S$ be a multiplicatively
	closed subset of $R$.
	A proper ideal $Q$ of $R$ is called an \emph{$S$-irreducible ideal} of $R$ if
	$S\cap Q=\varnothing$ and whenever
	$s(I\cap J)\subseteq Q\subseteq I\cap J$
	for some $s\in S$ and for some ideals $I,J$ of $R$,
	there exists $s'\in S$ such that
	$ss'I\subseteq Q$ or $ss'J\subseteq Q$.
\end{definition}

\begin{definition}
	Let $S$ be a multiplicatively closed subset of a $c$-lattice $L$. An element $q$ of $L$ such that $t\nleq q$ for all $t\in S$ is called \textit{$S$-irreducible element}, if $s(a\wedge b)\leq q \leq a\wedge b$ for some $s\in S$ and for some $a,b\in L$, then there exists an element $s^{\prime}\in S$ such that either $ss^{\prime}a\leq q$ or $ss^{\prime}b\leq q$.	
\end{definition}



\begin{lem}\label{thm:S-irreducible-ideal-lattice}
	Let $R$ be a commutative ring with identity and let $S$ be a multiplicatively
	closed subset of $R$.
	Let $L=Id(R)$ be the $c$-lattice of all ideals of $R$ and
	$S_L=\{(s)\in L \mid s\in S\}$.
	Then an ideal $Q$ of $R$ is an $S$-irreducible ideal of $R$ if and only if
	$Q$ (viewed as an element of $L$) is an $S_L$-irreducible element of $L$.
\end{lem}

\begin{proof}
	Assume that $Q$ is an $S$-irreducible ideal of $R$.
	Then $S\cap Q=\varnothing$.
	Hence $(s)\nleq Q$ for all $(s)\in S_L$.
	
	Let $A,B\in L$ and suppose that
	$(s)(A\wedge B)\leq Q\leq A\wedge B$
	for some $(s)\in S_L$.
	That is,
	$s(A\cap B)\subseteq Q\subseteq A\cap B$.

	As $Q$ is an $S$-irreducible ideal, there exists $s'\in S$ such that
	$ss'A\subseteq Q$ or $ss'B\subseteq Q$.
	Equivalently,
	$(ss')A\leq Q$ or $(ss')B\leq Q$ in $L$.
	Thus $Q$ is an $S_L$-irreducible element of $L$.
	
	Conversely, assume that $Q$ is an $S_L$-irreducible element of $L$.
	By the definition, $(s)\nleq Q$ for all $(s)\in S_L$.
	If $a\in S\cap Q$, then $(a)\leq Q$, a contradiction.
	Hence $S\cap Q=\varnothing$.
	
	Let $I,J$ be ideals of $R$ such that
	$s(I\cap J)\subseteq Q\subseteq I\cap J$
	for some $s\in S$.
	Equivalently,
	$(s)(I\wedge J)\leq Q\leq I\wedge J$ in $L$.
	Since $Q$ is $S_L$-irreducible, there exists $(s')\in S_L$ such that
	$(ss')I\leq Q$ or $(ss')J\leq Q$.
	That is,
	$ss'I\subseteq Q$ or $ss'J\subseteq Q$.
	Hence $Q$ is an $S$-irreducible ideal of $R$.
\end{proof}

\begin{prop}\label{pro1} 
	Let $L$ be an $S$-Noetherian lattice. Let $q_1,q_2,\dots,q_n\in L$, and $p$ be an $S$-prime element of $L$ such that $\displaystyle \bigwedge_{i=1}^{n} q_i\leq p$. Then there exists $s\in S$ such that $sq_k\leq p$ for some $k$. In particular, if $\displaystyle p=\bigwedge_{i=1}^{n} q_i$, then there exists $s\in S$ such that $sq_k\leq p\leq q_k$ for some $k$.
\end{prop}

\begin{proof} 
	Since $\displaystyle \prod_{i=1}^{n} q_i \leq \bigwedge_{i=1}^{n} q_i\leq p$ and $p$ is an $S$-prime element, there exists $s\in S$ such that $sq_k\leq p$ for some $k$. If $\displaystyle p=\bigwedge_{i=1}^{n} q_i$, then $sq_k\leq p\leq q_k$ for some $k$ and $s\in S$.
\end{proof}	



\begin{thm}\label{irred}
	Let $L$ be an $S$-Noetherian lattice. Then every $S$-irreducible element of $L$ is $S$-primary.
\end{thm}
\begin{proof}
	Let $q$ be a $S$-irreducible element of $L$. Let $a, m \in L$ such that $am\leq q$ with $sa\nleq q$ for all $s\in S$ and $m$ is a meet  principal element. Then consider the chain $((a\vee q): m)\leq ((a\vee q):m^2) \leq \dots \leq ((a\vee q):m^k) \leq \dots $. Since $L$ is an $S$-Noetherian lattice, by Theorem \ref{5.11a.}, there exists $s\in S$ and $k\in \mathbb{N}$ such that $s((a\vee q):m^{k+1})\leq ((a\vee q):m^k)$. 
	
	Let $c=a\vee q$ and $d=m^{k+1}\vee q$. Then $(c\wedge d):m^{k+1}=(c: m^{k+1})\wedge (d: m^{k+1})=((a\vee q):m^{k+1})\wedge((m^{k+1}\vee q):m^{k+1})=((a\vee q):m^{k+1})\wedge 1=(a\vee q):m^{k+1}$. 
	
	Therefore for any $s\in S$, $s((c\wedge d):m^{k+1})=s((a\vee q):m^{k+1})\leq (a\vee q):m^{k}$. Since $m^k$ and $m^{k+1}$ are meet principal elements, $((c\wedge d):m^{k+1})m^{k+1}=c\wedge d\wedge m^{k+1}$ and $((a\vee q):m^{k})m^k=(a\vee q)\wedge m^k$. 
	
	Clearly,  $ d\wedge (c\wedge d)=(q\vee m^{k+1})\wedge (c\wedge d)$. Since $L$ is an $r$-lattice, $L$ is  modular and  $q \leq c \wedge d $, hence by modularity, we have $(c\wedge d)\wedge(m^{k+1}\vee q)= (m^{k+1}\wedge c\wedge d)\vee q=q\vee ((c\wedge d):m^{k+1})m^{k+1}$, as $m^{k+1}$ is meet principal. Using the fact that $ x \leq (x : y ) y $, we have  $s(c \wedge d) \leq sq\vee s((c\wedge d):m^{k+1})m^{k+1}\leq q\vee ((a\vee q):m^{k})m^{k}m=q\vee ((a\vee q)\wedge m^k)m\leq q\vee (a\vee q)m=q\vee am\vee qm\leq q\vee q \vee q =q$. Hence $s(c\wedge d)\leq q$. Thus, we have $s(c\wedge d)\leq q\leq c\wedge d$. Since $q$ is a $S$-irreducible element, there exists $s^{\prime}\in S$ such that either $ss^{\prime}c\leq q$ or $ss^{\prime}d\leq q$. Since $ss^{\prime}c=ss^{\prime}a \vee ss^{\prime}q\nleq q$, as $sa\nleq q$ for all $s \in S$, we have $ss^{\prime}d\leq q$. Then $(ss^{\prime}m)^{k+1}\leq ss^{\prime}m^{k+1}\leq ss^{\prime}(m^{k+1}\vee q)= ss^{\prime}d\leq q$. This proves $ss^{\prime}m\leq \sqrt{q}$.
	
	Now, let $a,  b \in L$ such that  $ab\leq q$ with $sa\nleq q$ for all $s\in S$. Since $L$ is a $r$-lattice, $b$ is a finite join of meet principal elements, say, $b=m_1\vee m_2 \vee \dots \vee m_r$. Then $ab=am_1\vee am_2 \vee \dots \vee am_r\leq q$. This implies $am_i\leq q$ for all $i$. Since $sa\nleq q$ for all $s\in S$, there exists $s_i\in S$ such that $s_im_i\leq \sqrt{q}$ for $i=1,2,\dots ,r$. Take $s=s_1s_2\dots s_r$. Then $sb=sm_1\vee sm_2 \vee \dots \vee sm_r\leq \sqrt{q}$. This proves that $q$ is $S$-primary element of $L$.
\end{proof}

\begin{thm}[\textbf{Existence of $S$-Primary Decomposition}]\label{Existence} 
	Let $L$ be an $S$-Noetherian lattice. Then every proper element of $L$ can be written as a finite meet of $S$-primary elements.
\end{thm}
\begin{proof}
	Let $T$ be a collection of elements of $L$ that cannot be written as a finite meet of $S$-irreducible elements. Suppose $T\neq \emptyset$. Then by $S$-maximal condition, there exists an $S$-maximal element in $T$, say $m$. Then $m$ is not $S$-irreducible. Hence there exist $a,b \in L$ and $s \in S$ such that $	s(a \wedge b)\leq m \leq a \wedge b$,  and $\text{for all } s' \in S,\quad ss'a \nleq m \text{ and } ss'b \nleq m.$
	
	We claim that $a \notin T$. Suppose, on the contrary, that $a \in T$.  Since $m \leq a$ and $m$ is $S$-maximal in $T$, there exists $t \in S$ such that $ta \leq m$. Then  $(st)a \leq sm \leq m$.  Taking $s' = t$, we get $ss'a \leq m$,  a contradiction.
	
	Therefore $a\notin T$. Similarly $b\notin T$. This proves that $a$ and $b$ can be written as a finite meet of $S$-irreducible elements. Hence, $m$ can be written as a finite meet of $S$-irreducible elements. This gives $m\notin T$, a contradiction to $m\in T$. Therefore $T=\emptyset$. By Theorem \ref{irred}, every proper element of $L$ can be written as a finite meet of $S$-primary elements. 
\end{proof}


\begin{cor}\label{} 
	Let $L$ be an $S$-Noetherian lattice. Then every radical element $a$ of $L$ such that $t\nleq a$ for all $t\in S$ is the meet of finitely many $S$-prime elements.
\end{cor}
\begin{proof}
	By Theorem \ref{Existence}, there exists finitely many $S$-primary elements $q_1,q_2,\dots,q_n$ such that $a=q_1\wedge q_2\wedge \dots \wedge q_n$. Also $\sqrt{q_i}=p_i$ is $S$-prime for $i=1,2,\dots, n$. Hence $a=\sqrt{a}=\sqrt{q_1\wedge q_2\wedge \dots \wedge q_n}=\sqrt{q_1}\wedge \sqrt{q_2}\wedge \dots \sqrt{q_n}=p_1\wedge p_2\wedge \dots \wedge p_n$. 
\end{proof}




To prove the uniqueness theorem, we need the following definitions and terminologies. 


\begin{definition}
	Let $L$ be a multiplicative lattice with a multiplicatively closed subset $S \subseteq L_*$.  For any element $a \in L$, the \textit{$S$-saturation} of $a$ is $a_S = \bigvee \{ x \in L \mid sx \leq a \text{ for some } s \in S \}$.
	
	Note that in a $c$-lattice $L$, for any $x\in L_{\star}$, $x\leq a_S$ $\iff$ $sx\leq a$ for some $s\in S$.
	
	It is easy to verify that in a $c$-lattice, $a \leq a_S$ and $(a_S)_S = a_S$. Further, if $a \leq b$, then $a_S \leq b_S$.
\end{definition}

\begin{prop}\label{ps-prime}
	Let $L$ be an $c$-lattice and let $S \subseteq L$ be a multiplicatively closed subset. If $p \in L$ is an $S$-prime element, then its $S$-saturation $p_S$ is a prime element of $L$.
\end{prop}
\begin{proof}
	Let $a,b\in L_{\star}$ with $ab\leq p_S$. Then $sab\leq p$ for some $s\in S$. Since $p$ is $S$-prime, there exists $s'\in S$ such that $s'sa\leq p$ or $s'b\leq p$. Therefore $a\leq p_S$ or $b\leq p_S$. Thus $p_S$ is a prime element of $L$.
\end{proof}


 \begin{lem}\label{radical=ps} 
	Let $L$ be an $c$-lattice and $S \subseteq L$ be a multiplicatively closed subset. Let $q$ be an $S$-$p$-primary element of $L$ and $x\in L$ such that $sx\nleq q$ for all $s\in S$. Then $(q:sx)$ is $S$-primary element of $L$, and there exists $t\in S$ such that $t(\sqrt{(q:sx)})\leq p\leq \sqrt{(q:sx)}$ for all $s\in S$. In particular, if $sx \not\leq q$ for all $s \in S$ and $x\in L$, then
	$(\sqrt{q : x})_S = p_S$.
\end{lem}

\begin{proof} 
	First we prove that $(q:sx)$ is $S$-primary element for all $s\in S$. Let $s\in S$. If $s^{\prime}\leq (q:sx)$ for some $s^{\prime}\in S$, then $s^{\prime}sx\leq q$, a contradiction to $sx\nleq q$ for all $s\in S$. Therefore $s^{\prime}\nleq (q:sx)$ for all $s^{\prime}\in S$. 
	
	Let $a, ~b \in L$ such that  $ab\leq (q:sx)$ with $s^{\prime}a\nleq \sqrt{(q:sx)}$ for all $s^{\prime}\in S$. Then $(sxb)a\leq q$. Since $q$ is \linebreak $S$-primary element, there exists $s^{\prime\prime}\in S$ such that $s^{\prime\prime}sxb\leq q$ or $s^{\prime\prime}a\leq \sqrt{q}$. If $s^{\prime\prime}a\leq \sqrt{q}$, then $s^{\prime\prime}a\leq \sqrt{(q:sx)}$, a contradiction. Therefore $sxs^{\prime\prime}b \leq q$. Hence $s^{\prime\prime}b\leq (q:sx)$. Thus $(q:sx)$ is $S$-primary element of $L$. 
	
	Now, let $w$ be any compact element of $ L$ such that  $w\leq \sqrt{(q:sx)}$ for all $s \in S$. This implies $w^n\leq (q:sx)$ for some positive integer $n$. This implies $w^nsx\leq q$. Since $q$ is $S$-primary element and $sx\nleq q$ for all $s\in S$, there exists $t\in S$ such that $tw^n\leq \sqrt{q}$. Hence $(tw^n)^m\leq q$ for some positive integer $m$. This gives $(tw)^{mn}\leq (tw^n)^m\leq q$. Therefore $tw\leq \sqrt{q}$, i.e.,  $w \leq (\sqrt{q} : t)$. Since each compact element $\leq \sqrt{(q:sx)} $ is also $\leq (\sqrt{q} : t) $.  This gives $ \sqrt{(q:sx)} \leq (\sqrt{q} : t) $, i.e., $t(\sqrt{(q:sx)})\leq \sqrt{q}=p$. Also $q\leq (q:sx)$ implies $p=\sqrt{q}\leq \sqrt{(q:sx)}$. Thus, $t(\sqrt{(q:sx)})\leq p\leq \sqrt{(q:sx)}$ for all $s\in S$.	
	
	In particular, for $s=1$, $t(\sqrt{(q:x)})\leq p\leq \sqrt{(q:x)}$ for some $t\in S$. This implies $(\sqrt{q : x})_S = p_S$.
\end{proof}


\begin{thm}[First $S$-Uniqueness Theorem]
	Let $L$ be an $c$-lattice and let $S \subseteq L$ be a multiplicatively closed subset.  Let $a \in L$ be a minimal $S$-primary decomposition $ a=\bigwedge_{i=1}^{n} q_i,$
	where each $q_i$ is an \linebreak $S$-$p_i$-primary element of $L$ and $p_i = \sqrt{q_i}$. Then the set of $S$-saturated prime elements \linebreak $\{ (p_1)_S, (p_2)_S, \dots, (p_n)_S \}$ is uniquely determined by $a$ and is precisely the set \linebreak $	\mathbb{P} = \left\{\, (\sqrt{(a:x)})_S \;\middle|\; x \in L,\ sx \nleq a \ \text{for all } s \in S,\ \text{and } (\sqrt{(a:x)})_S \text{ is prime} \right\}$.
	
	In particular, the set $\{ (p_i)_S \}$ is independent of the chosen minimal $S$-primary decomposition of $a$.
\end{thm}
\begin{proof}
	We prove that the set of $S$-saturated primes $\{(p_1)_S,\dots,(p_n)_S\}$ coincides with $\mathbb{P}$.
	
	Fix $j \in \{1,\dots,n\}$. By minimality of the decomposition, $\bigwedge_{i \ne j} q_i \nleq (q_j)_S$. Hence there exists $x \in L$ such that $x \le \bigwedge_{i \ne j} q_i$ and $sx \nleq q_j$ for all $s \in S$. Since $a = \bigwedge_{i=1}^n q_i \le q_j$, it follows that $sx \nleq a$ for all $s \in S$.
	
	It is easy to observe that $(a:x) = \bigwedge_{i=1}^n (q_i:x)$. For $i \ne j$,  and as $x \le q_i$ implies $(q_i:x)=1$, hence $(a:x) = (q_j:x)$ and so we have $\sqrt{(a:x)} = \sqrt{(q_j:x)}$.
	
Since $sx \nleq q_j$ and $q_j$ is $S$-$p_j$-primary, we have, by Lemma \ref{radical=ps}, $(\sqrt{(q_j:x)})_S = (p_j)_S$. 
	
	Thus  $(\sqrt{(a:x)})_S = (p_j)_S$, showing $(p_j)_S \in \mathbb{P}$.
	
	Now, we prove the converse.
	
	Since $(a : x) = \left( \bigwedge_{i=1}^n q_i : x \right) = \bigwedge_{i=1}^n (q_i : x)$, we have $\sqrt{\bigwedge_{i=1}^n (q_i : x)} = \bigwedge_{i=1}^n \sqrt{(q_i : x)}$. This further gives $(\sqrt{(a:x)})_S =\left( \sqrt{\bigwedge_{i=1}^n (q_i : x)} \right)_S = \left( \bigwedge_{i=1}^n \sqrt{q_i : x} \right)_S= \bigwedge_{i=1}^n (\sqrt{q_i : x})_S $   as $a_S \wedge b_S = (a \wedge b)_S$. Hence $(\sqrt{(a:x)})_S \leq (\sqrt{q_i : x})_S$ for all $i$.
	
	Let $\pi = (\sqrt{(a:x)})_S \in \mathbb{P}$, where $sx \nleq a=\bigwedge_{i=1}^n q_i$ for all $s\in S$. 
	
	Since $\pi$ is prime and  $ \bigwedge_{i=1}^n (\sqrt{q_i : x})_S =(\sqrt{(a:x)})_S =\pi$,  there exists  some $j$ such that $(\sqrt{q_j : x})_S \leq (\sqrt{(a:x)})_S  $. This together with $(\sqrt{(a:x)})_S \leq (\sqrt{q_i : x})_S$ for all $i$ gives that $(\sqrt{(a:x)})_S = (\sqrt{q_j: x})_S$. 
	
	 If $sx \le q_j$ for some $s \in S$, then $s \le (q_j:x)$ and hence $(\sqrt{(q_j:x)})_S = 1$, contradicting the fact that $\pi$ is a proper prime element. Therefore $sx \nleq q_j$ for all $s \in S$, and by Lemma \ref{radical=ps}, $(\sqrt{(a:x)})_S =(\sqrt{(q_j:x)})_S = (p_j)_S$.\end{proof}
	 
	 
	

	\nocite{*}
\begin{thebibliography}{99}
\bibitem{HH} H. Ahmed and H. Sana, {\it Modules Satisfying the $S$-Noetherian property and $S$-ACCR}, Comm.  Algebra, {\bf 44 (5)} (2016), 1941-1951.
		
		\bibitem{AAJ} F. Alarcon, D. D. Anderson and C. Jayaram, {\it Some results on abstract commutative ideal theory}, Period. Math. Hung., {\bf 30} (1995), 1-26. 
		
		%\bibitem{ABT} F. A. A. Almahdi, E. M. Bouba and M. Tamekkante, {\it On weakly $s$-prime ideals of commutative rings}, Analele Stiint. ale Univ. Ovidius  Constanta Ser. Mat., Sciendo, {\bf 29 (2)} (2021), 173-186.
		
		\bibitem{A} D. D. Anderson, {\it Abstract commutative ideal theory without chain condition}, Algebra Universal., {\bf 6} (1976), 131-145. 
		
		\bibitem{AT} D. D. Anderson and T. Dumitrescu, {\it $S$-Noetherian Rings}, Comm. Algebra, {\bf 30 (9)} (2002), 4407-4416.
		
		
		
		%\bibitem{AB} D. D. Anderson and M. Bataineh, {\it Generalizations of prime ideals}, Commun. Algebra, 36 (2) (2008), 686-696. https://doi.org/10.1080/00927870701724177.
		
		
		%\bibitem{AB1} D. F. Anderson and A. Badawi, {\it On $n$-absorbing ideals of commutative rings}, Commun. Algebra, 39 (5) (2011), 1646-1672.  https://doi.org/10.1080/00927871003738998.
		
		%\bibitem{BGK} S. Ballal, M. Gophane and V. Kharat, {\it On weakly primary elements in multiplicative lattices}, Southeast Asian Bull. Math. {\bf 40 (1)} (2016), 49-57.
		
	    %\bibitem{B} R. G. Burton, {\it Fractional elements in multiplicative lattices}, Pacific J. Math. {\bf  56 (1)} (1975), 35-49.
		
		%\bibitem{C} G. C\u{a}lug\u{a}reanu, {\it Lattice concepts of module theory}, (Birkh\"{a}user, Basel, 1998).
		
		
		%\bibitem{CYT} F.  Callialp, E. Yetkin and U Tekir, {\it On $2$-absorbing primary and weakly $2$-absorbing elements in multiplicative lattices}, Ital. J. Pure Appl. Math., {\bf 34} (2015), 263-276.
		
		%\bibitem{CUU} E. Y. Celikel, E. A. Ugurlu  and G. Ulucak, {\it  On $\phi$-$2$-absorbing  elements in multiplicative lattices}, Palest. J. Math., 5 (1) (2016), 127-135.
		
		%\bibitem{CUU1} E. Y. Celikel, G. Ulucak and E. A. Ugurlu, {\it On $\phi$-$2$- absorbing primary elements in multiplicative lattices}, Palest. J. Math., 5 (1) (2016), 136-146.
		
		%\bibitem{CK} E. Celikel and H. Khashan, {\it On weakly $S$-primary ideals of a commutative rings}, J. Algebra Appl.,  https://doi.org/10.1142/S021949882450155X.
		
		
		\bibitem{D} R. P. Dilworth, {\it Abstract commutative ideal theory}, Pacific J. Math., {\bf 12} (1962), 481-498.
		
		%\bibitem{G} G. Gr\"atzer, {\it General Lattice Theory}, Birkh\"auser Basel (1998).
		
		\bibitem{HM} A. Hamed and A. Malek, {\it $S$-prime ideals of a commutative ring}, Beitr. Algebra Geom., {\bf 61} (2020),533-542.
		
		\bibitem{JS} V.  Joshi and S. Sarode, {\it Beck's conjecture and multiplicative
			lattices}, Discrete Math., {\bf 338 (3)} (2015), 93-98.
		
		
		%\bibitem{JB} V.  Joshi and S. Ballal, {\it A note on n-Baer multiplicative lattices}, Southeast Asian Bull. Math. {\bf 39} (2015), 67-76.
		
		%\bibitem{JS} V.  Joshi and S. Sarode, {\it Beck's conjecture and multiplicative lattices}, Discrete Math. {\bf 338 (3)}(2015), 93-98.
		
		
		%\bibitem{J} C. Jayaram, {\it Primary elements in Pr\"{u}fer lattices}, Czech. Math. J. {\bf 52 (127)} (2002), 585-593.
		
		
		\bibitem{JG} C. Jayaram, {\it Laskerian Lattices}, Czech. Math. J., {\bf 53 (128)} (2003), 351-363.
		
		%\bibitem{JJ} C. Jayaram and E. W. Johnson, {\it s-prime elements in multiplicative lattices}, Period. Math. Hung. {\bf 31} (1995), 201-208.
		
		%\bibitem{JJJ} C. Jayaram and  E. W. Johnson, {\it Strong compact elements in multiplicative lattices}, Czech. Math. J. {\bf 47 (122)} (1997), 105-112.
		
		%\bibitem{JTY} C. Jayaram, U. Tekir and E. Yetkin, {\it 2-absorbing and weakly 2-absorbing elements in multiplicative lattices}, Commun. Algebra, {\bf 42 (6)} (2014), 2338-2353. https://doi.org/10.1080/00927872.2012.761229.
		 
		 \bibitem{kap} I. Kaplansky,  \textit{Commutative Rings}, The University of Chicago Press: Chicago and London, 1974.
		 
	%	\bibitem{MMZ} N. Mahdou, M. A. S. Moutui and Y. Zahir, {\it On $S$-weakly prime ideals of commutative rings}, Georgian Mathematical Journal, {\bf 29 (3)} (2022), 397-405. https://doi.org/10.1515/gmj-2022-2141.
		
		
	%	\bibitem{MB} C. S. Manjarekar and A. V. Bingi, {\it  $\phi$-prime and $\phi$-primary elements in multiplicative lattices}, 	Algebra, Volume 2014, Article ID 890312, 7 pages.  https://doi.org/10.1155/2014/890312.
		
	%	\bibitem{MC} C. S. Manjarekar and N. S. Chavan, {\it Quasiregular and Baer lattices - I}, Period. Math. Hunger. {\bf 52} (2006), 31-49.
		
		%\bibitem{MC1} C. S. Manjarekar and A. N. Chavan, {\it $z$-elements and $z_j$-elements in multiplicative lattices}, Palestine J. Math. {\bf 8 (1)} (2019), 138-147.
		
		\bibitem{M} E. Massaoud, {\it $S$-primary ideals of a commutative ring}, Comm. Algebra, {\bf 50 (3)} (2021), 988–997. 
		
	%	\bibitem{MG} E. Massaoud and B. Gouaid, {\it Some results about weakly $S$-primary ideals of a commutative ring}, Gulf J. Math. {\bf 13 (1)} (2022), 88–100.
		
		%\bibitem{S} S. Sarode, {\it $\phi - \mathfrak{X}$ - elements in  multiplicative lattices}, J. Indian Math. Soc., {\bf 91 (1-2)} (2024), 217-228.
		
		\bibitem{SJ} S. Sarode and V. Joshi, {\it $\mathfrak {X} $-elements in multiplicative lattices-A generalization of $ J $-ideals, $ n $-ideals and $ r $-ideals in rings}, Int. Electron. J. Algebra, {\bf 32 (32)} (2022), 46-61.
		
		\bibitem{SJP} S. Sarode, V. Joshi and C. Patil, {\it On $S$-prime element principle} (	arXiv:2604.20820 ).
		
		\bibitem{SAK} T. Singh, A. Ansari and S. Kumar, {\it A study of $S$-primary decompositions}, Czech. Math. J., {\bf 75 (125)} (2025), 1241–1253.
		
		
		\bibitem{V} S. Visweswaran, {\it Some results on $S$-primary ideals of a commutative ring}, Beitr. Algebra Geom. {\bf 63} (2022), 247-266.
		
		
		
		\bibitem{WD} M. Ward and R. P. Dilworth, {\it Residuated lattices}, Trans. Amer. Math. Soc.,  45 (1939), 335-354.	
\end{thebibliography}

	
\end{document}





\bibitem{HH} Hamed Ahmed and Hizem Sana, {\it Modules Satisfying the $S$-Noetherian property and $S$-ACCR}, Communications in Algebra {\bf 44 (5)} (2016), 1941-1951.

\bibitem{AAJ} F. Alarcon, D. D. Anderson and C. Jayaram, {\it Some results on abstract commutative ideal theory}, Period. Math. Hung. {\bf 30 (1)} (1995), 1-26.

\bibitem{AD} D. D. Anderson, {\it Abstract commutative ideal theory without chain condition}, Algebra Universalis {\bf 6} (1976), 131-145.

\bibitem{AT} D. D. Aderson and Tiberiu Dumitrescu, {\it $S$-Noetherian Rings}, Commun. Algebra {\bf 30 (9)} (2002), 4407-4416.


%\bibitem{AJ} D. D. Anderson and E. W. Johnson, {\it Join-principally generated multiplicative lattices}, Algebra Universalis {\bf 19} (1984), 74-82.

\bibitem{BGK} S. Ballal, M. Gophane and V. Kharat, {\it On weakly primary elements in multiplicative lattices}, Southeast Asian Bull. Math. {\bf 40 (1)} (2016), 49-57.

\bibitem{B} R. G. Burton, {\it Fractional elements in multiplicative lattices}, Pacific J. Math. {\bf  56 (1)} (1975), 35-49.

\bibitem{C} G. C\u{a}lug\u{a}reanu, {\it Lattice concepts of module theory}, (Birkh\"{a}user, Basel, 1998).

\bibitem{Dil} R. P. Dilworth, {\it Abstract commutative ideal theory}, Pacific J. Math. {\bf 12} (1962), 481-498.





%\bibitem{FR} L. Fuchs and R. Reis, {\it On lattice-ordered commutative semigroups}, Algebra Universalis {\bf 50 (3)} (2003), 341-357.

\bibitem{G} G. Gr\"atzer, {\it General Lattice Theory}, Birkh\"auser Basel (1998).




\bibitem{HM} A. Hamed and A. Malek, {\it $S$-prime ideals of a commutative ring}, Beitr Algebra Geom {\bf 61} (2020),533-542.




\bibitem{JB} V.  Joshi and S. Ballal, {\it A note on n-Baer multiplicative lattices}, Southeast Asian Bull. Math. {\bf 39} (2015), 67-76.

\bibitem{JS} V.  Joshi and S. Sarode, {\it Beck's conjecture and multiplicative
	lattices}, Discrete Math. {\bf 338 (3)}(2015), 93-98.


\bibitem{J} C. Jayaram, {\it Primary elements in Pr\"{u}fer lattices}, Czech. Math. J. {\bf 52 (127)} (2002), 585-593.


\bibitem{JG} C. Jayaram, {\it Laskerian Lattices}, Czech. Math. J. {\bf 53 (128)} (2003), 351-363.

\bibitem{JJ} C. Jayaram and E. W. Johnson, {\it s-prime elements in multiplicative lattices},
Period. Math. Hung. {\bf 31} (1995), 201-208.

\bibitem{JJJ} C. Jayaram and  E. W. Johnson, {\it Strong compact elements in multiplicative lattices}, Czech. Math. J. {\bf 47 (122)} (1997), 105-112.



\bibitem{JTY} C. Jayaram, U. Tekir and E. Yetkin, {\it 2-absorbing and weakly 2-absorbing elements in multiplicative lattices}, Communications in Algebra {\bf 42} (2014), 1-16.

%\bibitem{JTY1} C. Jayaram, U. Tekir,  E. Yetkin, {\it On 2-absorbing primary and weakly 2-absorbing  elements in multiplicative lattices}, Italian j. pure and appl. math. {\bf 34} (2015), 263-276.


\bibitem{KJ} S. Kavishwar and V. Joshi, {\it A note on prime ideal principle in lattices}, Asian-European J. Math. {\bf 8 (2)} (2015) 1550024 (12 pages).


%\bibitem{KB} H. A. Khashan and A. B. Bani-Ata, {\it $J$-ideals of Commutative rings}, IEJA. 

\bibitem{lr} T. Y. Lam  and M. L. Reyes, {\it A prime ideal principle in commutative algebra}, J. Alebra. {\bf 319  }(2008), 3006-3027.
%\bibitem[2008]{lr1} Lam. T. Y and Reyes. M. L, {\it Oka and Ako ideal families in commutative rings},



\bibitem{MB} C. S. Manjarekar and A. V. Bingi, {\it $\phi$-prime and $\phi$-primary elements in multiplicative lattices}, 
Hindawi Publishing Corporation
Algebra, Volume 2014, Article ID 890312, 7 pages http://dx.doi.org/10.1155/2014/890312

\bibitem{MC} C. S. Manjarekar and N. S. Chavan, {\it Quasiregular and Baer lattices - I}, Period. Math. Hunger. {\bf 52} (2006), 31-49.

\bibitem{MC1} C. S. Manjarekar and A. N. Chavan, {\it $z$-elements and $z_j$-elements in multiplicative lattices}, Palestine J. Math. {\bf 8 (1)} (2019), 138-147.


%\bibitem{M} R. Mohamadian, {\it $r$-ideals in commutative rings}, Turk. J. Math. {\bf 39} (2015), 733-749.

%\bibitem{TKO} U. Tekir, S. Koc and K. H. Oral, {\it $n$-Ideals of commutative rings}, Filomat {\bf 31 (10)} (2017), 2933-2941.

\bibitem{SJ} S. Sarode and V. Joshi, {\it $\mathfrak {X} $-elements in multiplicative lattices-A generalization of $ J $-ideals, $ n $-ideals and $ r $-ideals in rings}, Int. Electron. J. Algebra, DOI:10.24330/ieja.1102289 (also see  https://arxiv.org/abs/2101.06667).







%\bibitem{man} N. K. Thakare and  C. S.  Manjarekar, {\it Abstract spectral theory: multiplicative lattices in which every character is contained in a unique maximal character}, Algebra and its Applications,  Lecture Notes in Pure and Appl. Math., Dekker, New York, {\bf 91} (1984), 265-276.

%\bibitem{TMM} N. K. Thakare, C. S. Manjarekar and S. Maeda, {\it Abstract Spectral theory II: Minimal characters and minimal spectrums of multiplicative lattices }, Acta Sci. Math. (Szeged)
%{\bf 52} (1988), 53-67.



\bibitem{SS} S. Sarode, {\it $\phi - \mathfrak{X}$ - elements in  multiplicative lattices}.

\bibitem{SJP} S. Sarode, V. Joshi and C. Patil, {\it On $S$-prime element principle} (Under preparation).


\bibitem{wd} M. Ward and R. P. Dilworth, {\it Residuated lattices}, Trans. Amer. Math. Soc. {\bf 45} (1939), 335-354.






 Recently, Hamed and Malek \cite{HM} defined the concept of $S$-prime ideals in commutative rings, which is a generalization of prime ideals in commutative rings with identity. An ideal $P$ of commutative ring $R$ with identity is called an $S$-prime ideal of $R$ if there exists an $s \in S$ such that for all $a, b \in R$ if $ab \in P$, then $sa \in P$ or $sb \in P$, where $S$ is a multiplicative subset of $R$ disjoint with $P$. Analogously, Sarode, joshi and Patil \cite{SJP} introduced the $S$-prime elements in $V$-lattices. Let  $S$ be a multiplicatively closed set of a $V$-lattice $L$. A proper element $p$ of $L$ such that $t \not\leqq p$ for all $t \in S$. Then $p$ is an $S$-prime element of $L$ if there exists an element $s \in S$ such that for all $a, b  \in L$ with $a \cdot b \leq p$, implies $s \cdot a \leq p$ or $s \cdot b \leq p$.  Massaoud \cite{M} and  Visweswaran \cite{V} introduced independently the concept of an $S$-prime ideals which is a generalization of prime ideals  in a commutative ring $R$ with identity.  Also, the concept of weakly $S$-prime ideals as a generalization of weakly prime ideals is independently introduced by  Celikel and Khashan \cite{CK} and Massaoud and Gouaid \cite{MG} in a commutative ring $R$ with identity. In this paper, we introduce  the concept of  $S$-prime  and weakly $S$-prime elements which is generalizations of  prime  and weakly prime elements in multiplicative lattices. Let $S$ be a multiplicatively closed subset of a $V$-lattice $L$. A proper element $p$ of $L$ such that $t\not\leqq p$ for all $t\in S$ is called an $S$-prime (weakly $S$-prime) if there exists an element $s\in S$ such that for all $a,b\in L$ with $ab\leq p$ ($(ab \neq 0)\leq p$), implies $sa\leq p$ or $sb\leq \sqrt p$.  
 We prove that an $S$-prime ideal (weakly $S$-prime ideal) of a commutative ring $R$ with identity, where $S$ is a multiplicatively closed subset of $R$, is a $S_L$-prime element (weakly $S_L$-prime element) of a $V$-lattice $L = Id(R)$ with respect to some multiplicatively closed subset $S_L$ of $L$.





If every element of multiplicative lattice $L$ is a compact element, then $L$ is called \textit{compact} lattice. If $1$ is a compact element of $L$, then $L$ is called \textit{$1$-compact} lattice. By 	 a {\it $V$-lattice}, we  mean a  multiplicative lattice which is
1-compact, compactly generated in which the
product of two compact elements is compact. A nonempty subset $S$ of $L_{*}$ in a $V$-lattice $L$ is \textit{multiplicatively closed} if $a \cdot b \in S$ whenever $a, b \in S$.	Throughout this paper, we consider a multiplicatively closed subset $S$ of a $V$-lattice $L$ such that $1\in S$ and $0\notin S$. Let $L$ be a multiplicative lattice and $a,b\in L$. Then  $(a:b) =
\bigvee \{x \;|\;  x\cdot b \leq a\}$. 	
Let $L$ be a $V$-lattice and $a\in L$. The the radical of $a$ is denoted by $\sqrt a$, is defined as $\sqrt a =
\bigvee \{x\in L_* \;|\;  x^n \leq a ~ \textrm{for some}~ n\in \mathbb{Z}^+\}$. 




\begin{lem}\label{} 
	Let $S$ be a multiplicatively closed subset of a $V$-lattice $L$ and $r$ is a proper element of $L$ such that $r\in S$. If $p$ is an $S$-prime element of $L$, then $pr$ is  an $S$-prime element of $L$.
\end{lem}

\begin{proof} Let  $r$ be a proper element of $L$ such that $r\in S$ and  $p$ is an $S$-prime element of $L$. Since $pr\leq p$ and  $t \not\leqq p$ for all $t \in S$, we get $t\not\leqq pr$ for all $t\in S$. Let $c,d\in L$ such that $cd\leq pr$. Therefore, $cd\leq p$. Since $p$ is an $S$-prime element of $L$, there exists an element $s\in S$ such that $sc\leq p$ or $sd\leq p$. Hence, $sc\leq p$ or $s^nd^n\leq p$ for some positive integer $n$. This implies $scr\leq pr$ or $s^nd^nr^n\leq pr$. Since $s,r\in S$, let $s^{\prime}=sr\in S$. Therefore, $s^{\prime}c\leq pr$ or $s^{\prime}d\leq \sqrt{pr}$. Thus, $pr$ is an $S$-prime element of $L$.
\end{proof}	



\begin{definition}
	Let $S$ be a multiplicatively closed subset of a $V$-lattice $L$. A proper element $p$ of $L$ such that $t\not\leqq p$ for all $t\in S$ is called $S$-prime if there exists an element $s\in S$ such that for all $a,b\in L_{\star}$ with $ab\leq p$, $sa\leq p$ or $sb\leq \sqrt p$.	
\end{definition}


%\begin{lem}\label{cc} Let $S$ be a multiplicatively closed subset of a $V$-lattice $L$. Let $p$ be a proper element of $L$ such that $t\not\leqq p$ for all $t\in S$. Then the following are equivalent:\begin{enumerate}
\item\label{1.9.1} $p$ is an $S$-prime element of $L$.
\item\label{1.9.2} There exists an element $s\in S$ such that for all $a,b\in L_{\star}$ with $ab\leq p$, implies $sa\leq p$ or $sb\leq \sqrt p$.
\end{enumerate}
\end{lem}

\begin{proof} 
	(\ref{1.9.1}) $\implies$ (\ref{1.9.2}) obvious.
	
	(\ref{1.9.2}) $\implies$ (\ref{1.9.1}) Let $a,b\in L$ with $ab\leq p$ and $sa \not\leqq p$ for all $s \in S$. As $L$ is a $V$-lattice and $sa \not\leqq p$ for all $s \in S$, There exists $(c \neq 0) \in L_*$ such that $c \leq sa$ for all $s \in S$ and $c \not\leqq p$. Now let any $d \in L_*$ such that $d \leq sb$ for all $s \in S$. Since $c \cdot d \leq sa \cdot sb \leq a\cdot b \leq p$, by hypothesis there exists an element $s\in S$ such that $sd \leq p$. Thus every compact element $\leq sb$ is also $\leq p$. We get $sb \leq p$. Hence $p$ is an $S$-prime element of $L$.
\end{proof}

	Firstly, we show that $p= \sqrt{{\wedge^{\infty}_{n=1} p_n}}= \wedge^{\infty}_{n=1} \sqrt{p_n}$. Since $p = \bigwedge^{\infty}_{n=1} p_n \leq p_n$ for all $n \in \mathbb{Z}^{+}$, we get $p \leq \sqrt{p_n}$ for all  $n \in \mathbb{Z}^{+}$ and hence $p \leq \wedge^{\infty}_{n=1} \sqrt{p_n}$. Let $x \in L_*$ such that $x \leq \wedge^{\infty}_{n=1} \sqrt{p_n} \leq \sqrt{p_n}$ for all  $n \in \mathbb{Z}^{+}$. By Lemma \ref{3.4},  $ \wedge^{\infty}_{n=1} \sqrt{p_n} =\sqrt{0}$. Therefore $x^m = 0$ for some positive integer $m$. So $x^m \leq p_n$ for all  $n \in \mathbb{Z}^{+}$. This gives $x^m \leq \wedge^{\infty}_{n=1} \sqrt{p_n} = p $ and hence $x \leq p$. Therefore every compact element $\leq \wedge^{\infty}_{n=1} \sqrt{p_n}$ is $\leq p$. Hence  $p= \sqrt{{\wedge^{\infty}_{n=1} p_n}}= \wedge^{\infty}_{n=1} \sqrt{p_n}$.
	
	
	
	
	
	
	
	{\color{red} \begin{proof}If $L$ is an $S$-Noetherian lattice, then every element is $S$-compact and hence every prime element is $S$-compact. 
			
			
			Conversely, assume that every $S$-prime element of $L$  is $S$-compact. Let $\beta = \{a' \in L\; |\; \ t \not\leqq a'$ for all $t \in S$ and $a'$ is not  $S$-compact$\}$. If $\beta = \emptyset$, then all elements of $L$ are $S$-compact. By Theorem \ref{5.11a.}, $L$ is an $S$-Noetherian lattice.
			
			Now, assume that $\beta \neq \emptyset$. Let $a_{\alpha} \in \beta$, $\alpha \in \Lambda$ ($\Lambda$ is an index set). Let $ a_1 \leq a_2 \leq a_3 \leq \cdots$ be increasing sequence of elements from $\beta$. Let $a = \bigvee_{\alpha} a_{\alpha}$, $\alpha \in \Lambda$ ($\Lambda$ is an index set). We show that $a$ is upper bound of $a_1 \leq a_2 \leq a_3 \leq \cdots$ and $a \in \beta$.
			
			First, we claim that $a$ is not an $S$-compact element of $L$. Suppose on the contrary that $a$ is an $S$-compact element of $L$. Therefore there exists $s \in S$ and a compact element $b \in L$ such that $s \cdot  a \leq b \leq a$, that is $s \cdot (\bigvee_{\alpha} a_{\alpha}) \leq b \leq (\bigvee_{\alpha} a_{\alpha})$. Since $b$ is compact and  $b \leq (\bigvee_{\alpha} a_{\alpha})$, we have $b \leq a_1 \vee a_2 \vee \cdots \vee a_n$ for some $ n \in \mathbb{Z}^{+}$. As  $a_1 \leq a_2 \leq a_3 \leq \cdots$ is ascending chain of elements and $b \leq a_1 \vee a_2 \vee \cdots \vee a_n=a_n$ for some $n \in \mathbb{N}$,  we have $s \cdot (\bigvee_{\alpha} a_{\alpha})  \leq b \leq a_n$. Since $s \cdot a_n \leq s \cdot (\bigvee_{\alpha} a_{\alpha})$, we have $s \cdot  a_n \leq b\leq a_n$, a contradiction to $a_n$ is not an $S$-compact element. Therefore  $a$ is not an $S$-compact element of $L$. Further, observe that $t \nleq a$ for all $t \in S$. If possible $t_1 \leq a$ for some $t_1 \in S$, we get $\displaystyle t_1\leq \bigvee_{\alpha} a_{\alpha}$. Since $t_1$ is compact, $\displaystyle t_1\leq \bigvee_{i=1}^{n} a_{\alpha_i}=a_{\alpha_j}$ for some $a_{\alpha_j}$; $1\leq j\leq n$, a contradiction to $t\nleq a_{i}$ for all $t \in S$. Hence,  $a \in \beta$. Moreover, $a$ is an upper bound of  $a_1 \leq a_2 \leq a_3 \leq \cdots$. By Zorn's Lemma, $\beta$ contains a maximal element, say $p$. If $p$ is prime, then $p$ is also $S$-prime for the multiplicatively closed subset $S$ of $L$. Hence by the hypothesis, $p$ is an $S$-compact element, contradicting $p\in\beta$ and in this case we are through. 
			
			If $p$ is not a prime element, then, by Remark \ref{princ-prime}, there exist principal elements $x, y$ in $L$ such that $xy \leq p$ with $x \not\leqq p$ and $y \not\leqq p$. Let $A = x \vee p$ and $B = (p : x)$. Clearly, $p<A,B$, as $x\leq A$, but $x\nleq p$ and $y\leq B$, but $y\nleq p$. Therefore, $A,B\notin \beta$. If $s\leq A=p\vee x$ for some $s\in S$, then $sy\leq py\vee xy\leq p$. Thus, $p$ is an $S$-prime element. Hence, by hypothesis, $p$ is $S$-compact, a contradiction to $p\in \beta$. Thus, $t\nleq A$ for all $t\in S$. We can also show that $t\nleq B$ for all $t\in S$. Therefore, $A$ and $B$ are $S$-compact elements of $L$. Since $x$ is principal, $p \wedge x = x \cdot (p: x) = x \cdot  B$. By Lemma \ref{5.51.}, $x \cdot B$ is $S$-compact and hence $p \wedge x$ is $S$-compact. Since $L$ is modular, by  isomorphism Theorem for modular lattices, we have 
			$[p\wedge x, p] \cong [x, p\vee x]$ and hence $p = y \vee (p \wedge x)$, where $y$ is a compact element. By Lemma \ref{5.51.}, $p$ is $S$-compact, contradicting $p\in\beta$.
			Thus, $\beta=\varnothing$, i.e., every element of $L$ is $S$-compact, and by Theorem~\ref{5.11a.}, the lattice $L$ is $S$-Noetherian.
	\end{proof}}
